% Root proof; B independent complete literal audit passed for both statements.
\subsection{A projective central corner in a star}

\begin{theorem}\label{C:star-projective-center22}
Let $C$ have outward star support with at least two leaves and
at most twenty-two arrows. Suppose the trivial left module of
its central endomorphism monoid is projective over its rational
monoid algebra. Then $BC\not\simeq\mathbb CP^2$.
For an inward star the same conclusion holds with the trivial
right module of the central monoid in place of the left module.
\end{theorem}
\begin{proof}
Assume $BC\simeq\mathbb CP^2$. Let $M$ be the central monoid,
let $N_i$ be the leaf monoids, and let ${}_{N_i}S_i{}_M$ be
the nonempty cross bisets. Apply the initial-object version
of Proposition~\ref{C:finite-target-extreme-corner} at the center.
It gives an actual functor $Q:C\to\bar C$ replacing $M$ by
the trivial monoid and $S_i$ by the left $N_i$-set
$T_i=S_i/M$. Write $Y=B\bar C$.
The map $BQ$ is onto on fundamental groups, induces rational
and $\mathbb F_{65521}$ cohomology-ring isomorphisms, and has
a homotopy left inverse $h:Y\to BC$. The prime is good
because every relevant finite group has order at most $22$.

There is no nonidentity arrow between different branches.
Thus $Y$ is the actual wedge of the ordinary left-act cones
\[
 Y_i=B\mathcal C^{\rm left}(N_i,T_i)
\]
at their common trivial central vertex. Each $Y_i$ is connected.
Since $Y$ is simply connected, the free-product formula implies
that every $Y_i$ is simply connected. The finite category nerve
has finitely many simplices in each degree, so all its homology
groups in a fixed degree are finitely generated. In particular,
simple connectivity and rational second rank one give
$H^2(Y;\mathbb Z)\cong\mathbb Z$.

Let $u$ generate $H^2(BC;\mathbb Z)$ and put $v=h^*u$.
The equation $h\,BQ\simeq1$ shows that $v$ is an integral
generator, and that $v^2$ is nonzero and primitive modulo
torsion in $H^4(Y;\mathbb Z)$. The direct sum decomposition
of second cohomology for this finite wedge puts $v$ in one
unique summand $Y_j$; all other integral second-cohomology
groups are zero. Write $v_j$ for its restriction there.
The rational cohomology of $Y$ is that of the projective plane,
and $v_j^2\ne0$, so $Y_j$ has exactly that rational homology.
The square $v_j^2$ is primitive modulo torsion: divisibility
in this wedge summand would give the same divisibility of
$v^2$ in the direct sum.

The map $h$ restricted to $Y_j$, followed by a chosen equivalence
$BC\simeq\mathbb CP^2$, pulls the projective-plane generator
back to $v_j$. Miller's theorem therefore makes the restriction
of $v_j$ zero along every map from a finite group. Over
$\mathbb F_{65521}$, the class $v^2$ remains nonzero because
$(BQ)^*v^2=u^2$. It lies in the same wedge summand $Y_j$.
As the total Betti numbers in degrees two, three and four
are $(1,0,1)$, those of $Y_j$ must also be $(1,0,1)$.

This ordinary act cone has at most twenty arrows. Indeed,
there is at least one other nonempty branch, whose leaf
identity and one cross arrow contribute two arrows; hence
\[
 |N_j|+|T_j|+1
 \leq |C|-|M|-2+1\leq20.
\]
All four conditions of
Theorem~\ref{C:finite-target-small-act-twenty} now hold for
$Y_j$, after opposition to turn its left act into a right act.
This is impossible. Taking opposites proves the inward version.
No finite-CW-type assertion about $Y$ or its wedge summands
has been used.
\end{proof}

\begin{corollary}\label{C:star-contractible-center18}
A finite inward or outward star with at least two leaves,
at most eighteen arrows, and contractible central monoid
classifying space cannot have type $\mathbb CP^2$.
In particular this applies if the central monoid has a
left zero or a right zero.
\end{corollary}
\begin{proof}
The full branch nerves meet in the contractible central
nerve, so their actual union is homotopy equivalent to
their wedge. The free-product fundamental-group formula
and the integral cohomology ring of $\mathbb CP^2$ then
force one branch to have type $\mathbb CP^2$: precisely one
summand supports both the degree-two generator and its
nonzero square, and every other simply connected summand
is acyclic and contractible. Each branch has at most
sixteen arrows, since a different branch contributes at
least two. Theorem~\ref{C:CP2-lower17} excludes it.
A left or right zero gives a natural transformation
between the identity and a constant functor on the
one-object category, proving the stated sufficient condition.
\end{proof}

\begin{corollary}\label{C:star-small-center18}
An inward or outward star category with at least two leaves,
at most eighteen arrows, and central endomorphism monoid of
order at most four cannot have type $\mathbb CP^2$.
\end{corollary}
\begin{proof}
Take opposites to make the star outward and write $M$ for
its center. If the trivial left rational $M$-module is
projective, the preceding theorem applies. Otherwise the
minimal ideal has at least two rows. Its identity lies
outside that ideal. Since $|M|\leq4$, its minimal Rees
data must have one column and trivial maximal group.
A minimal idempotent $e$ then satisfies $em=e$ for every
$m\in M$. It is a left zero, so $BM$ is contractible.

The nerve of the star is the union of its branch nerves,
intersecting in the full central nerve $BM$. These are
simplicial cofibrations. Their homotopy pushout over the
contractible common subspace is therefore a wedge of the
individual two-object branch spaces. If this wedge has
type $\mathbb CP^2$, one branch has that type. Indeed each
branch is simply connected by the free-product formula;
reduced integral homology splits over the wedge, and the
primitive nonzero square lies in its unique degree-two
summand. That same summand carries degree four. All other
summands are simply connected and acyclic, hence contractible,
and the chosen branch is homotopy equivalent to the wedge.

Every branch has at most sixteen arrows, since another
branch contributes at least its leaf identity and one
cross arrow. Theorem~\ref{C:CP2-lower17} excludes it.
\end{proof}
